\documentclass[a4paper,11pt,dvipsnames]{article}
\usepackage{amsmath,amsthm,amssymb,amsfonts}
\usepackage{array}
\usepackage{tabularx}
\usepackage[table]{xcolor}
\usepackage[most]{tcolorbox}
\usepackage{graphicx}
\graphicspath{{./img/}}
\usepackage[a4paper,left=1.5cm,right=1.5cm,top=1cm,bottom=1.5cm]{geometry}
% \usepackage[hmargin=1cm,vmargin=1cm]{geometry}
\usepackage{array}
\usepackage{setspace}
\newcommand*\diff{\mathop{}\!\mathrm{d}}
\newcommand*\Diff[1]{\mathop{}\!\mathrm{d^#1}}

\newcommand{\N}{\mathbb{N}}
\newcommand{\Z}{\mathbb{Z}}

% A style for tcolorbox without frame etc. "notcb"
\tcbset{notcb/.style={enhanced jigsaw, colback=white, coltitle=black,sharp corners, boxrule=0pt}}

% Style for the number tables
\tcbset{tabularheader/.style={arc=2pt,notcb,boxsep=0pt,left=0pt,colbacktitle=white,rounded corners=all,enhanced jigsaw,coltitle=black,tabularx={*{6}{Y|}Y},boxrule=1pt,attach boxed title to top left, boxed title style={enhanced jigsaw,sharp corners,left=0pt,boxrule=0pt}}}

\newcolumntype{Y}{>{\centering\arraybackslash}X}

\usepackage[locale=FR, per-mode=fraction, separate-uncertainty=true]{siunitx}
\usepackage{physics}
\usepackage[utf8]{inputenc}
\usepackage[T1]{fontenc}
\usepackage[spanish]{babel}

\usepackage{enumitem}

\usepackage{pgfplots}
% \usepackage{pdfpages}
\usepackage{adjustbox, tikz}
% \usepackage{gnuplottex}
\usepackage{xcolor}
\usepackage[locale=FR, per-mode=fraction, separate-uncertainty=true]{siunitx}
\usepackage{physics}

\usepackage{ifpdf}
\ifpdf%
  \usepackage{pdftricks}
  \begin{psinputs}
    \usepackage{pst-optic}
%\usepackage{pstricks-add} 	
  \end{psinputs}
\else
  \usepackage{pst-optic}
  \newenvironment{pdfpic}{}{}
\fi
 

\usetikzlibrary{decorations}
\usetikzlibrary{decorations.pathmorphing}
\usetikzlibrary{shapes.geometric}
\usetikzlibrary{arrows}
\usetikzlibrary{arrows.meta}   %edg
\usetikzlibrary{calc}
\pgfplotsset{compat=newest}


\begin{document}

{\Large Resoluciones 1er parcial - 2022}

\begin{enumerate}[leftmargin=*]
  \item Un cilindro de $\SI{50.0}{\centi\meter^3}$ de cobre está inicialmente a $\SI{20.0}{\celsius}$. ¿A qué temperatura su volumen aumentará $\SI{0.075}{\centi\meter^3}$? Coeficiente de expansión lineal del cobre: $\alpha = \SI{1.7E-5}{\celsius^{-1}}$.\\
  \begin{tabularx}{\textwidth}{XXXXXX}
    $\square$ $\SI{14.7}{\celsius}$ & $\square$ $\SI{29.4}{\celsius}$ & $\square$ $\SI{34.7}{\celsius}$ & $\square$ $\SI{49.4}{\celsius}$ & $\square$ $\SI{64.1}{\celsius}$ & $\square$ $\SI{108.2}{\celsius}$\\
  \end{tabularx}

\

\textit{Resolución:}
\begin{flalign*}
    \Delta V &= \beta V_0 \Delta T
\end{flalign*}

El coeficiente de expansión volumétrica es: $\beta = 3\alpha$.
\begin{flalign*}
    \SI{0.075}{\centi\cubic\meter}  &= 3\cdot \SI{1.7E-5}{\celsius^{-1}}\cdot \SI{50.0}{\centi\cubic\meter} \cdot \Delta T\\
    \Delta T  &=  \SI{29.4}{\celsius} = T_f - \SI{20.0}{\celsius}
\end{flalign*}

Respuesta: $T = \SI{49.4}{\celsius}$

\

  \item Un refrigerador tiene un rendimiento de $2.10$. Durante cada ciclo, absorbe $\SI{3.40E4}{\joule}$ de calor del depósito frío. ¿Cuánta energía mecánica se requiere en cada ciclo para operar el refrigerador?\\
  \begin{tabularx}{\textwidth}{XXXXXX}
    $\square$ $\SI{0.86E4}{\joule}$ & $\square$ $\SI{1.31E4}{\joule}$ & $\square$ $\SI{1.62E4}{\joule}$ & $\square$ $\SI{3.40E4}{\joule}$ & $\square$ $\SI{5.02E4}{\joule}$ & $\square$ $\SI{5.84E4}{\joule}$\\
  \end{tabularx}

  \

\textit{Resolución:}
\begin{flalign*}
    e &= \frac{Q_f}{W}\\
    2.10 &= \frac{\SI{3.40E4}{\joule}}{W}
\end{flalign*}

Respuesta: $W = \SI{1.62E4}{\joule}$


  \item Un recipiente aislado térmicamente y de capacidad calorífica despreciable contiene $\SI{0.285}{\kilogram}$ de un metal en estado sólido a $\SI{530}{\celsius}$. En dicho recipiente se introduce $\SI{1.00}{\kilogram}$ de una mezcla de agua y hielo que se encontraba inicialmente en equilibrio, y el sistema alcanza el equilibrio térmico a $\SI{4.3}{\celsius}$. \textit{Datos}: calor específico del agua líquida: $\SI{4190}{\joule.\kilogram^{-1}.\kelvin^{-1}}$; calor específico del hielo: $\SI{2100}{\joule.\kilogram^{-1}.\kelvin^{-1}}$; calor específico del metal: $\SI{410}{\joule.\kilogram^{-1}.\kelvin^{-1}}$; calor de fusión del agua: $\SI{334E3}{\joule.\kilogram^{-1}}$.\\
  a) Encontrar la masa de hielo contenida en la mezcla que se agrega al recipiente.\\
  b) Calcular la variación de energía interna del metal.\\
  c) Calcular la variación de entropía del metal. 

  \

\textit{Resolución:}\\
\textit{a})
\begin{flalign*}
    m_\text{metal}C_\text{metal}(T_\text{eq}-T_\text{i,metal}) + m_\text{hielo}L_f + m_\text{agua}C_\text{agua}(T_\text{eq}-\SI{0}{\celsius}) &= 0\\
    \SI{0.285}{\kilogram}\cdot\SI{410}{\joule.\kilogram^{-1}.\kelvin^{-1}}\cdot(\SI{4.3}{\celsius}-\SI{530}{\celsius}) + m_\text{hielo}\cdot\SI{334E3}{\joule.\kilogram^{-1}}&\\
     + \SI{1.00}{\kilogram}\cdot\SI{4190}{\joule.\kilogram^{-1}.\kelvin^{-1}}\cdot(\SI{4.3}{\celsius}-\SI{0}{\celsius}) &=0
\end{flalign*}

Respuesta: $m_\text{hielo} = \SI{0.13}{\kilogram}$

\

\textit{b})
\begin{flalign*}
    \Delta U_\text{metal} &= Q_\text{metal} = \SI{0.285}{\kilogram}\cdot\SI{410}{\joule.\kilogram^{-1}.\kelvin^{-1}}\cdot(\SI{4.3}{\celsius}-\SI{530}{\celsius})
\end{flalign*}

Respuesta: $\Delta U_\text{metal} = \SI{-61.4}{\kilo\joule}$

\

\textit{c})
\begin{flalign*}
    \Delta S_\text{metal} &= mC\ln(T_f/T_i) = \SI{0.285}{\kilogram}\cdot\SI{410}{\joule.\kilogram^{-1}.\kelvin^{-1}}\cdot\ln(\SI{277.3}{\kelvin}/\SI{803}{\kelvin})
\end{flalign*}

Respuesta: $\Delta S_\text{metal} = \SI{-124.2}{\joule/\kelvin}$

\

  \item Una pared de $\SI{14}{\square\meter}$ que está compuesta por $\SI{1.5}{\centi\meter}$ de un material aislante y $\SI{15}{\centi\meter}$ de ladrillo, separa un ambiente con aire a $\SI{25}{\celsius}$ en contacto con el material aislante, del exterior a una temperatura de $\SI{-4}{\celsius}$. Los coeficientes de conductividad térmica del material aislante y del ladrillo son $\SI{0.04}{\watt.\meter^{-1}.\kelvin^{-1}}$ y $\SI{0.6}{\watt.\meter^{-1}.\kelvin^{-1}}$ respectivamente. Se sabe que la potencia calórica transmitida por esta pared es $\SI{550}{\watt}$ mientras las temperaturas del aire interior y exterior permanecen constantes. Asumiendo que los coeficientes de convección para el aire son iguales a ambos lados de la pared, encuentre la temperatura sobre la superficie libre del ladrillo.\\

\

\textit{Resolución:}\\
Primero es necesario encontrar el valor del coeficiente de convección.
\begin{flalign*}
    H &= \frac{1}{\frac{2}{h}+\frac{e_1}{k_1}+\frac{e_2}{k_2}}A\Delta T\\
    \SI{550}{\watt} &= \frac{1}{\frac{2}{h}+\frac{\SI{0.015}{\meter}}{\SI{0.04}{\watt.\meter^{-1}.\kelvin^{-1}}}+\frac{\SI{0.15}{\meter}}{\SI{0.6}{\watt.\meter^{-1}.\kelvin^{-1}}}}\SI{14}{\square\meter} \cdot (\SI{25}{\celsius}-(\SI{-4}{\celsius}))\\
    h &= \SI{17.67}{\watt.\meter^{-2}.\kelvin^{-1}}
\end{flalign*}

Variación de temperatura desde la superficie exterior de la pared hasta la temperatura del aire exterior:
\begin{flalign*}
    H &= hA\Delta T\\
    \SI{550}{\watt} &= \SI{17.67}{\watt.\meter^{-2}.\kelvin^{-1}}\cdot \SI{14}{\square\meter} \cdot (T-(\SI{-4}{\celsius}))
\end{flalign*}
Respuesta: $T = \SI{-1.78}{\celsius}$

\

  \begin{minipage}[t]{.5\textwidth}
  \item El ciclo de la figura aproxima el funcionamiento de una máquina térmica constituida por $\SI{0.55}{\mole}$ de un gas ideal diatómico que intercambia calor con dos depósitos que mantienen sus temperaturas constantes. El gas comienza en el estado $a$ con un volumen de $\SI{2.3}{\liter}$ y a la temperatura del reservorio caliente, $\SI{520}{\kelvin}$. Se expande adiabáticamente hasta un volumen de $\SI{9.0}{\liter}$, llegando al estado $b$ cuya temperatura es la del reservorio frío, $\SI{300}{\kelvin}$. Luego se comprime isotérmicamente hasta el estado $c$ con un volumen de $\SI{1.5}{\liter}$. Por último evoluciona sobre la trayectoria rectilínea que une los estados $c$ y $a$ para completar el ciclo. Calcular la eficiencia térmica de esta máquina. 
    \end{minipage}
  \hspace{0.02\linewidth}
  \begin{minipage}[t]{.4\textwidth}
    \strut\vspace*{-\baselineskip}\newline
  \begin{center}
    \begin{tikzpicture}[scale=0.9]
      \begin{axis}[
        ticks=none,
        axis x line=bottom,
        axis y line=left,
        xmin=0, xmax=10, 
        ymin=0, ymax=12.5, 
        xlabel={$V$},
        ylabel={$p$},
        % xtick={2,9},
        % xticklabels={1.5,9},
        % ytick={1.5},
        % yticklabels={1.5}
        %yticklabels={},
        %extra x ticks={0.04},
        %extra y ticks={2}
        ];
      \addplot [-latex, color=blue, very thick] [samples= 180, domain=2.3:4.5]  {1.5*(9/x)^1.5};
      \addplot [color=blue, very thick] [samples= 180, domain=4.5:9]  {1.5*(9/x)^1.5};
      \addplot [latex-, color=blue, very thick] [samples= 180, domain=3:9]  {1.5*9/x};
      \addplot [color=blue, very thick] [samples= 180, domain=1.5:3]  {1.5*9/x};
      % \addplot[color = black, dashed, thick] coordinates {(2, 0) (2, 1.5) (0, 1.5)};
      % \addplot[color = black, dashed, thick] coordinates {(9, 0) (9, 1.5)};
      \draw [color=blue, very thick][-latex](1.5,9)--(2,10.63);
      \draw [color=blue, very thick](2,10.63)--(2.3,11.61);
      \draw (2.3,11.61) node [right] {$a$};
      \draw (9,1.5) node [above right] {$b$};
      \draw (1.5,9) node [left] {$c$};
      \fill [black](2.3,11.61) circle(2pt);
      \fill [black](9,1.5) circle(2pt);
      \fill [black](1.5,9) circle(2pt);
      \end{axis}
    \end{tikzpicture}
  \end{center}
\end{minipage}

\

\textit{Resolución:}\\
Rendimiento de una máquina: $e=W/Q_\text{absorbido}$. El calor absorbido es la suma de todos los calores positivos durante el ciclo. En este caso, el único calor positivo es el intercambiado en el proceso $c\rightarrow a$. Como no tenemos una fórmula para calcular ese calor, se puede encontrar conociendo la variación de energía interna y el trabajo.

\begin{flalign*}
  \Delta U_{ca} &= nC_v(T_a-T_c) = \SI{0.55}{\mol}\cdot \frac{5}{2} \cdot\SI{8.31}{\joule\mol^{-1}\kelvin^{-1}}(\SI{520}{\kelvin}-\SI{300}{\kelvin})\\
  \Delta U_{ca} &= \SI{2514}{\joule}
\end{flalign*}

El trabajo en el proceso $c\rightarrow a$ se puede calcular como el área debajo de la curva:
\begin{flalign*}
  W_{ca} &= (V_a-V_c)p_c+(V_a-V_c)(p_a-p_c)/2
\end{flalign*}
Las presiones en $a$ y $c$ se pueden encontrar con la ecuación de estado del gas ideal:
\begin{flalign*}
  p_a &= \frac{nRT_a}{V_a} = \frac{\SI{0.55}{\mol}\cdot \SI{8.31}{\joule\mol^{-1}\kelvin^{-1}}\cdot\SI{520}{\kelvin}}{\SI{0.0023}{\cubic\meter}} = \SI{1.03E6}{\pascal}\\
  p_c &= \frac{nRT_c}{V_c} = \frac{\SI{0.55}{\mol}\cdot \SI{8.31}{\joule\mol^{-1}\kelvin^{-1}}\cdot\SI{300}{\kelvin}}{\SI{0.0015}{\cubic\meter}} = \SI{0.914E6}{\pascal}
\end{flalign*}
Entonces el trabajo en ese proceso es:
\begin{flalign*}
  W_{ca} = (\SI{0.0023}{\cubic\meter}-\SI{0.0015}{\cubic\meter})\cdot \SI{0.914E6}{\pascal}+&\\(\SI{0.0023}{\cubic\meter}-\SI{0.0015}{\cubic\meter})(\SI{1.03E6}{\pascal}-\SI{0.914E6}{\pascal})/2 = \SI{777.6}{\joule}
\end{flalign*}
Ahora se puede obtener el calor en $c\rightarrow a$, que es igual al calor absorbido en el ciclo:
\begin{flalign*}
  Q_{ca} &= \Delta U_{ca} + W_{ca} = \SI{2514}{\joule} + \SI{777.6}{\joule} = \SI{3292}{\joule}
\end{flalign*}
Para obtener el trabajo del ciclo se necesitan los trabajos de cada proceso, falta calcular los trabajos en $a\rightarrow b$ y $b\rightarrow c$:
\begin{flalign*}
  W_{ab} &= -\Delta U_{ab} = -nC_v(T_b-T_a) = \SI{2514}{\joule}\\
  W_{bc} &= nRT_b\ln(V_c/V_b) = \SI{0.55}{\mol}\cdot \SI{8.31}{\joule\mol^{-1}\kelvin^{-1}}\cdot\SI{300}{\kelvin}\cdot\ln(\SI{0.0015}{\cubic\meter}/\SI{0.009}{\cubic\meter}) = \SI{-2457}{\joule}
\end{flalign*}
El trabajo del ciclo resulta:
\begin{flalign*}
  W &= W_{ab} + W_{bc} + W_{ca}\\
  W &= \SI{2514}{\joule} - \SI{2457}{\joule} + \SI{777.6}{\joule}\\
  W &= \SI{834.6}{\joule}
\end{flalign*}

Respuesta:
\begin{flalign*}
  e = \frac{\SI{834.6}{\joule}}{\SI{3292}{\joule}} = 0.2535
\end{flalign*}


\end{enumerate}

\end{document}
